\documentclass[10pt]{article}
% Surcouche compacte pour les interros
% Socle commun dérivé du framework des khôlles SI
% Version autonome et compatible pdfLaTeX pour le dépôt Interros.
\RequirePackage[T1]{fontenc}
\RequirePackage[utf8]{inputenc}
\RequirePackage[french]{babel}
\RequirePackage{lmodern}
\RequirePackage[margin=1.8cm]{geometry}
\RequirePackage{xcolor}
\RequirePackage{graphicx}
\RequirePackage{amsmath,amssymb}
\RequirePackage{tikz}
\RequirePackage{tabularx}
\RequirePackage{array}
\RequirePackage{enumitem}
\RequirePackage[most]{tcolorbox}
\RequirePackage{fancyhdr}
\RequirePackage{hyperref}
\RequirePackage{titlesec}
\usetikzlibrary{arrows.meta,positioning,calc,fit,chains,decorations.pathmorphing,decorations.markings,shapes,shapes.multipart,patterns,graphs,quotes}
\definecolor{SIblue}{RGB}{31,78,121}
\definecolor{SIlightblue}{RGB}{229,240,250}
\definecolor{SIred}{RGB}{190,0,0}
\definecolor{SIgray}{RGB}{235,235,235}
\hypersetup{colorlinks=true,linkcolor=SIblue,urlcolor=SIblue}
\pagestyle{fancy}
\setlength{\headheight}{14pt}
\fancyhf{}
\lhead{Sciences industrielles pour l'ingénieur}
\rhead{Interros ATS}
\cfoot{\thepage}
\titleformat{\section}{\Large\bfseries\color{SIblue}}{}{0pt}{}
\titleformat{\subsection}{\large\bfseries\color{SIblue}}{}{0pt}{}
\setlist[itemize]{leftmargin=*,itemsep=2pt,topsep=2pt}
\setlist[enumerate]{leftmargin=*,itemsep=4pt,topsep=4pt}
\newtcolorbox{donnees}[1][]{enhanced,breakable,colback=white,colframe=SIblue,title=Données,fonttitle=\bfseries,#1}
\newtcolorbox{attention}[1][]{enhanced,breakable,colback=orange!8,colframe=orange!80!black,title=Attention,fonttitle=\bfseries,#1}
\newtcolorbox{methode}[1][]{enhanced,breakable,colback=green!5,colframe=green!45!black,title=Méthode,fonttitle=\bfseries,#1}
\newcommand{\torseur}[2]{\left\{\mathcal V_{#1}\right\}_{#2}}
% Bibliothèques SI simplifiées, compatibles avec les interros
\providecommand{\og}{``}
\providecommand{\fg}{''}

% Graphe de liaisons
\tikzset{
  glsolide/.style={circle,draw=SIblue,very thick,minimum size=10mm,fill=SIlightblue,font=\bfseries},
  glliaison/.style={draw=SIred,very thick,rounded corners=2pt,font=\small,align=center}
}
\newenvironment{grapheliaisons}[1][]{\begin{center}\begin{tikzpicture}[node distance=25mm,#1]}{\end{tikzpicture}\end{center}}
\newcommand{\GLsolide}[3]{\node[glsolide] (#1) at (#3) {#2};}
\newcommand{\GLliaison}[4][]{\draw[glliaison,#1] (#2) -- node[midway,fill=white,inner sep=2pt]{#4} (#3);}
\newcommand{\GLpivot}[4][]{\GLliaison[#1]{#2}{#3}{Pivot\\$#4$}}
\newcommand{\GLglissiere}[4][]{\GLliaison[#1]{#2}{#3}{Glissière\\$#4$}}
\newcommand{\GLencastrement}[3][]{\GLliaison[#1]{#2}{#3}{Encastrement}}
\newcommand{\GLhelicoidale}[4][]{\GLliaison[#1]{#2}{#3}{Hélicoïdale\\$#4$}}

% Schémas cinématiques simplifiés
\newenvironment{schemacinematique}[1][]{\begin{center}\begin{tikzpicture}[line cap=round,line join=round,#1]}{\end{tikzpicture}\end{center}}
\newcommand{\SCbati}[2]{\coordinate (#1) at (#2);\draw[thick] ($(#1)+(-.65,0)$)--($(#1)+(.65,0)$);\foreach \x in {-.55,-.3,...,.55}{\draw ($(#1)+(\x,0)$)--++(-.18,-.18);}}
\newcommand{\SCpivot}[2]{\node[draw,circle,minimum size=5mm,thick] (#1) at (#2) {};}
\newcommand{\SCglissiere}[2]{\node[draw,rectangle,minimum width=10mm,minimum height=5mm,thick] (#1) at (#2) {};\draw[thick] ($(#1)+(-.35,-.18)$)--($(#1)+(.35,.18)$);}
\newcommand{\SChelicoidale}[2]{\node[draw,circle,minimum size=5mm,thick] (#1) at (#2) {};\draw[thick] ($(#1)+(-.15,0)$) arc (180:-90:.15);}
\newcommand{\SCrelier}[3][]{\draw[very thick,#1] (#2)--(#3);}

% SysML simplifié
\newenvironment{diagrammeSysML}[1][]{\begin{center}\begin{tikzpicture}[node distance=18mm,#1]}{\end{tikzpicture}\end{center}}
\newcommand{\SYSacteur}[4][]{\node[draw,circle,minimum size=3mm,#1] (head#4) at (#2) {};\draw (head#4)--++(0,-.55) coordinate (body#4);\draw ($(body#4)+(0,.35)$)--++(-.28,-.22);\draw ($(body#4)+(0,.35)$)--++(.28,-.22);\draw (body#4)--++(-.25,-.35);\draw (body#4)--++(.25,-.35);\node[below=8mm of head#4] (#4) {#3};}
\newcommand{\SYScas}[4][]{\node[draw,ellipse,align=center,minimum width=28mm,minimum height=10mm,#1] (#4) at (#2) {#3};}
\newcommand{\SYSinclude}[2]{\draw[dashed,->] (#1)--node[above]{\scriptsize\og include\fg}(#2);}
\newcommand{\SYSextend}[2]{\draw[dashed,->] (#1)--node[above]{\scriptsize\og extend\fg}(#2);}

% Tableau de mobilités
\newcommand{\mobilites}[6]{\begin{tabular}{c|c}R&T\\\hline #1&#2\\#3&#4\\#5&#6\end{tabular}}


\geometry{a4paper,left=13mm,right=13mm,top=18mm,bottom=16mm,headheight=28pt,headsep=8mm,footskip=10mm}
\fancyhf{}
\renewcommand{\headrulewidth}{0pt}
\renewcommand{\footrulewidth}{0.6pt}
\fancyfoot[L]{\footnotesize Lycée Robert Doisneau}
\fancyfoot[R]{\footnotesize \thepage}
\newcommand{\InterroHeader}[2]{%
\noindent\begin{tcolorbox}[enhanced,boxrule=.7pt,arc=2mm,colback=SIblue,colframe=SIblue,left=3mm,right=3mm,top=1.5mm,bottom=1.5mm]
\begin{minipage}[c]{.18\linewidth}\color{white}\bfseries\Large SI\end{minipage}%
\begin{minipage}[c]{.60\linewidth}\centering\color{white}\bfseries\large #1\end{minipage}%
\begin{minipage}[c]{.20\linewidth}\raggedleft\color{white}\bfseries Interro\\de cours\end{minipage}
\end{tcolorbox}
\vspace{-1mm}
\noindent\begin{minipage}[c]{.78\linewidth}\begin{tcolorbox}[height=12mm,colback=white,colframe=black,boxrule=1pt,arc=1.5mm,left=2mm,top=2mm]\textbf{Nom :}\end{tcolorbox}\end{minipage}\hfill
\begin{minipage}[c]{.13\linewidth}\begin{tcolorbox}[height=12mm,colback=white,colframe=black,boxrule=1pt,arc=2mm,halign=center,valign=center]\textbf{#2}\end{tcolorbox}\end{minipage}
\vspace{2mm}}
\newcounter{interroquestion}
\newcommand{\question}[2]{\refstepcounter{interroquestion}\par\medskip\noindent\textbf{Question \theinterroquestion :}\quad #1\hfill\textbf{(#2 points)}\par\smallskip}
\newtcolorbox{reponse}[1][35mm]{colback=white,colframe=black,boxrule=.5pt,arc=0mm,height=#1}
\newcommand{\resetquestions}{\setcounter{interroquestion}{0}}
\newcommand{\chaineenergie}{\begin{tikzpicture}[font=\small,>=Latex,node distance=4mm,bloc/.style={draw,minimum width=25mm,minimum height=10mm,align=center}]
\node[bloc](a){\dots};\node[bloc,right=of a](b){MODULER};\node[bloc,right=of b](c){\dots};\node[bloc,right=of c](d){\dots};\node[bloc,fill=yellow!30,right=of d](e){AGIR};
\draw[->,thick](a)--(b);\draw[->,thick](b)--(c);\draw[->,thick](c)--(d);\draw[->,thick](d)--(e);\end{tikzpicture}}
\newcommand{\chaineinformation}{\begin{tikzpicture}[font=\small,>=Latex,node distance=5mm,bloc/.style={draw,minimum width=25mm,minimum height=11mm,fill=yellow!30,align=center}]
\node[bloc](a){\dots};\node[bloc,right=of a](b){\dots};\node[bloc,right=of b](c){\dots};\draw[->,very thick,SIblue](a)--(b);\draw[->,very thick,SIblue](b)--(c);\end{tikzpicture}}

\usepackage{array}
\setlength{\parindent}{0pt}
\renewcommand{\arraystretch}{1.25}

\newcommand{\choixliaisons}{%
\begin{tabularx}{\linewidth}{|X|X|X|X|X|}\hline
A) Pivot & B) Pivot glissant & C) Glissière & D) Sphère-plan & E) Sphérique\\\hline
\end{tabularx}}

\newcommand{\symboleSpherique}{%
\begin{tikzpicture}[baseline=-1mm,scale=.65]
\draw[SIblue,very thick] (-.8,0)--(-.25,0);
\draw[SIblue,very thick] (0,0) circle (.27);
\fill[SIred] (0,0) circle (.07);
\end{tikzpicture}}

\newcommand{\symbolePivot}{%
\begin{tikzpicture}[baseline=-1mm,scale=.65]
\draw[SIblue,very thick] (-1,0)--(-.35,0);
\draw[SIblue,very thick] (0,0) circle (.38);
\draw[SIred,very thick] (0,0) circle (.25);
\draw[SIred,very thick] (.25,0)--(1,0);
\node[SIblue] at (-.72,.26) {1};\node[SIred] at (.65,.3) {2};
\end{tikzpicture}}

\begin{document}
\InterroHeader{Modélisation cinématique}{10}

\question{À quelle liaison correspond ce schéma : \symboleSpherique}{0,5}
\choixliaisons

\question{À quelle liaison correspond ce schéma : \symbolePivot}{0,5}
\choixliaisons

\question{Quelle est la bonne caractéristique d'une liaison pivot ?}{0,5}
\begin{tabularx}{\linewidth}{|X|X|X|X|X|}\hline
A) de centre $A$ & B) de direction $\vec{x}$ & C) de direction $(A,\vec{x})$ & D) d'axe $(A,\vec{x})$ & E) de normale $(A,\vec{x})$\\\hline
\end{tabularx}

\question{Quelle liaison possède deux degrés de liberté en translation et trois degrés de liberté en rotation ?}{0,5}
\choixliaisons

\medskip
Soit une base $(\vec{x}_1,\vec{y}_1,\vec{z}_1)$ en rotation par rapport à la base $(\vec{x}_0,\vec{y}_0,\vec{z}_0)$. L'angle $\theta$ caractérise la rotation autour de l'axe $\vec{z}_0=\vec{z}_1$ : $\theta=(\vec{x}_0,\vec{x}_1)$.

\smallskip
Soit une base $(\vec{x}_2,\vec{y}_2,\vec{z}_2)$ en rotation par rapport à la base $(\vec{x}_1,\vec{y}_1,\vec{z}_1)$. L'angle $\beta$ caractérise la rotation autour de l'axe $\vec{y}_1=\vec{y}_2$ : $\beta=(\vec{z}_1,\vec{z}_2)$.

\question{Dessiner les deux figures de changement de base.}{2}
\begin{reponse}[46mm]\end{reponse}

\question{Pour chacun des trois mécanismes, déterminer la liaison équivalente des liaisons montées en parallèle.}{3}
\begin{center}
\begin{tabular}{|p{.31\linewidth}|p{.31\linewidth}|p{.31\linewidth}|}\hline
\centering
\begin{tikzpicture}[scale=.72,>=Latex]
\draw[->] (0,0)--(2.4,0) node[right]{$\vec{x}$};
\draw[->] (0,0)--(0,1.8) node[above]{$\vec{z}$};
\draw[->] (0,0)--(1.5,1.15) node[right]{$\vec{y}$};
\draw[SIred,very thick,rounded corners=10mm] (-.2,-.4)--(2.8,2.1)--(3.4,1.5)--(.5,-.9)--cycle;
\foreach \p/\n in {(1.0,.55)/A,(2.15,1.35)/B}{\draw[SIblue,very thick] \p circle(.18);\node[above left] at \p {$\n$};}
\end{tikzpicture}
&\centering
\begin{tikzpicture}[scale=.72,>=Latex]
\draw[SIblue,very thick] (-1.4,-.5)--(1.2,-.5)--(1.2,.8)--(.2,.8);
\draw[green!60!black,very thick] (-.9,.2)--(.4,.2)--(.4,1.2)--(-.9,1.2)--cycle;
\draw[green!60!black,very thick] (-.9,.2) circle(.18);\node[above left] at (-.9,.2){$B$};
\draw[green!60!black,very thick] (.4,.2) circle(.18);\node[below right] at (.4,.2){$A$};
\node[green!60!black] at (-.2,.75){1};\node[SIblue] at (1.35,.55){0};
\draw[->] (-1.6,-.9)--(-.7,-.9)node[right]{$\vec{x}$};
\draw[->] (-1.6,-.9)--(-1.6,.05)node[above]{$\vec{z}$};
\end{tikzpicture}
&\centering
\begin{tikzpicture}[scale=.72,>=Latex]
\draw[->] (-.2,0)--(3.4,0) node[right]{$\vec{y}$};
\draw[->] (1.5,-.7)--(1.5,1.5) node[above]{$\vec{z}$};
\draw[SIred,very thick] (.3,0)--(2.8,0);
\draw[SIred,very thick] (.3,0) circle(.25);\node[above] at (.3,.3){$A$};
\draw[SIred,very thick] (2.8,0) circle(.25);\node[above] at (2.8,.3){$B$};
\draw[SIblue,very thick] (-.1,-.4)--(.7,-.4);\foreach \x in {0,.2,...,.6}{\draw[SIblue] (\x,-.4)--++(-.15,-.2);}
\draw[SIblue,very thick] (2.4,-.4)--(3.2,-.4);\foreach \x in {2.5,2.7,...,3.1}{\draw[SIblue] (\x,-.4)--++(-.15,-.2);}
\node at (1.5,-.25){$O$};
\end{tikzpicture}
\tabularnewline[2mm]\hline
\rule{0pt}{31mm} & & \\\hline
\end{tabular}
\end{center}

\newpage
\InterroHeader{Modélisation cinématique}{10}

Soit le robot manipulateur suivant :
\begin{center}
\begin{tikzpicture}[scale=.9,>=Latex]
\draw[very thick] (-3,-1.6)--(-3,1.4);\foreach \y in {-1.5,-1.25,...,-.8}{\draw (-3.35,\y)--(-3,\y+.2);}
\draw[SIblue,very thick] (-3,.3) circle(.25);\node[left] at (-3,.3){$A$};
\draw[SIblue,very thick] (-2.75,.3)--(-.7,-.2);
\draw[SIblue,very thick] (-.7,-.2) circle(.24);\node[below] at (-.7,-.2){$B$};
\draw[SIred,very thick] (-.45,-.2)--(1.45,-.55);
\draw[SIred,very thick] (1.45,-.55) rectangle (1.85,.15);
\draw[brown!70!black,very thick] (1.65,-1.5)--(1.65,1.35);
\draw[->] (-3,.3)--(-1.2,1.4)node[above]{$\vec{y}_2$};
\draw[->] (-3,.3)--(.3,1.55)node[above]{$\vec{y}_1$};
\draw[->] (-3,.3)--(-1.5,-1.3)node[below]{$\vec{x}_1$};
\draw[->] (-.7,-.2)--(.1,.9)node[above]{$\vec{z}_3$};
\draw[->] (1.65,-.2)--(2.8,-.2)node[right]{$\vec{x}_3$};
\draw[->] (1.65,-.2)--(2.75,.65)node[right]{$\vec{x}_4$};
\draw[->] (-3,.3)--(-3,2)node[above]{$\vec{z}_1=\vec{z}_2$};
\node[font=\bfseries] at (-2.7,-.65){1};\node[font=\bfseries] at (-1.1,.35){2};\node[font=\bfseries] at (.8,.1){3};\node[font=\bfseries] at (1.25,.8){4};
\node at (-1.3,-.9){$\theta_{21}$};\node at (1.2,-.95){$\theta_{32}$};\node at (2.15,.45){$\lambda_{43}$};
\end{tikzpicture}
\end{center}

\question{Dessiner les figures de changement de base associées aux paramètres $\theta_{21}$ et $\theta_{32}$.}{1}
\begin{reponse}[48mm]\end{reponse}

\question{Donner le graphe de liaisons du robot manipulateur.}{2}
\begin{reponse}[93mm]\end{reponse}

\end{document}
